\chapter{Implementierung des VirtIO-GPU-Treibers}

Die Implementierung des VirtIO-GPU-Treibers in hhuOS ist eine mehrschichtige
softwaretechnische Aufgabe. Sie überbrückt die Kluft zwischen hochrangigen
Grafikaufrufen und der Kommunikation mit der zugrunde liegenden PCI-Hardware.
Der Code ist architektonisch in mehrere Submodule unter
\texttt{src/device/graphic/virtio/} gegliedert. Diese Aufteilung stellt sicher,
dass Discovery, Initialisierung, Virtqueue-Transport, Control-Queue-Kommandos,
Ressourcenverwaltung und Laufzeitpfade klar voneinander getrennt sind.

\section{Geräteerkennung und Probe-Mechanismus}
Der Lebenszyklus des Treibers beginnt mit der Erkennung der Hardware. In der
hhuOS-Umgebung wird dies von der Funktion
\texttt{virtio\_gpu\_probe\_all\_devices()} verwaltet.

\subsection{PCI-Enumeration und Identifikation}
Während des Kernel-Bootvorgangs, insbesondere im Einstiegspunkt
\texttt{GatesOfHell.cpp} oder über einen Autostart-Hook, wird der PCI-Bus
durchsucht. Der Treiber sucht nach der Hersteller-ID \texttt{0x1AF4}
(Standard für VirtIO) und der Geräte-ID \texttt{0x1050} (VirtIO-GPU). Im
Gegensatz zu älteren Treibern, die Speicheradressen fest kodieren, liest diese
Implementierung dynamisch die Base Address Registers (BARs), um den
Konfigurationsraum des Geräts zu lokalisieren.

\subsection{Capability-Durchlauf und Abbildung}
Ein wesentlicher Bestandteil der Implementierung ist die Routine
\texttt{build\_regions\_from\_caps()}. VirtIO-1.2-Geräte nutzen eine verkettete
Liste von PCI-Capabilities, um ihre Registerbereiche zu definieren
\cite{Virtio12}. Der Treiber durchläuft diese Liste, um die obligatorischen
Regionen \texttt{COMMON\_CFG}, \texttt{NOTIFY\_CFG} und \texttt{ISR\_CFG} sowie
optional \texttt{DEVICE\_CFG} zu identifizieren.

Die einzelnen Bereiche erfüllen dabei folgende Aufgaben:
\begin{itemize}
    \item \textbf{Common Configuration:} Wird für die Statusverwaltung, die Feature-Verhandlung und die Auswahl bzw. Konfiguration der Virtqueues verwendet.
    \item \textbf{Notification Region:} Dient als „Türklingel“ (Doorbell), um dem Gerät neue Einträge im Available Ring zu signalisieren.
    \item \textbf{ISR Configuration:} Dient der Rückmeldung von Interrupt- beziehungsweise Completion-relevanten Ereignissen.
    \item \textbf{Device Configuration:} Stellt optionale gerätespezifische Daten wie Display-Parameter bereit.
\end{itemize}

Für jeden gefundenen Bereich ruft der Treiber die HAL-Funktion
\texttt{os\_map\_mmio()} auf, um eine virtuelle Abbildung auf die physischen
Hardware-Register zu erzeugen und diese für die weitere Initialisierung im
Kernel verfügbar zu machen.

\begin{center}
\begin{tikzpicture}[node distance=1cm, >=Stealth, block/.style={rectangle, draw, fill=blue!10, text width=4cm, text centered, minimum height=1cm}]
    \node [block] (pci) {PCI-Scan (\texttt{pci\_find\_first})};
    \node [block, right=of pci] (caps) {Capability-Liste durchlaufen};
    \node [block, right=of caps] (mmio) {BAR-Bereiche abbilden (MMIO)};
    \node [block, below=of mmio] (init) {\texttt{virtio\_gpu\_try\_init()}};

    \draw [->, thick] (pci) -- (caps);
    \draw [->, thick] (caps) -- (mmio);
    \draw [->, thick] (mmio) -- (init);

    \node[draw, dashed, fit=(pci) (caps) (mmio), label=above:{\textbf{Probe-Phase (Kernel-Integration)}}] {};
\end{tikzpicture}
\captionof{figure}{Detaillierte Treibererkennungs- und Abbildungspipeline}
\end{center}

\section{Der VirtIO-Handshake und die Initialisierung}
Nach der Abbildung der MMIO-Bereiche tritt der Treiber in die Funktion
\texttt{virtio\_gpu\_try\_init()} ein. Diese folgt der strikten
Zustandsmaschine der VirtIO-Spezifikation und stellt sicher, dass weder
Queue-Zugriffe noch GPU-Kommandos stattfinden, bevor der Host die
ausgehandelten Eigenschaften akzeptiert hat.

\subsection{Feature-Verhandlung}
In \texttt{negotiate\_features()} liest der Treiber die vom Gerät angebotenen
Features als 64-Bit-Bitmaske in zwei 32-Bit-Fenstern aus. Anschließend wird
konservativ nur die tatsächlich benötigte Teilmenge aktiviert. Zwingend ist
dabei \texttt{VIRTIO\_F\_VERSION\_1}, da nur damit der moderne VirtIO-Pfad mit
Little-Endian-Datenstrukturen und speicherabgebildeten Capabilities genutzt
werden kann.

Die Feature-Verhandlung folgt dabei dem üblichen Ablauf:
\begin{enumerate}
    \item Lesen der vom Gerät angebotenen Feature-Bits,
    \item Auswahl der vom Treiber unterstützten Teilmenge,
    \item Schreiben der akzeptierten Teilmenge in die Treiberregister,
    \item Setzen von \texttt{FEATURES\_OK},
    \item erneute Validierung, dass das Gerät dieses Bit weiterhin gesetzt lässt.
\end{enumerate}

Zusätzliche Features wie VirGL werden erkannt und intern vermerkt, sind jedoch
nicht Bestandteil des aktuellen 2D-Hauptpfads. Die Implementierung verfolgt
damit bewusst eine konservative Politik: Stabilität und Kompatibilität haben
Vorrang vor einer frühzeitigen Aktivierung optionaler Erweiterungen.

\subsection{Statusübergänge}
Die Implementierung verwendet das Gerätestatusregister, um den Fortschritt der
Initialisierung explizit zu signalisieren:
\begin{enumerate}
    \item \textbf{ACKNOWLEDGE:} Der Treiber hat das Gerät erkannt.
    \item \textbf{DRIVER:} Der Treiber signalisiert, dass er das Gerät grundsätzlich ansteuern kann.
    \item \textbf{FEATURES\_OK:} Gast und Host stimmen dem ausgehandelten Feature-Satz zu.
    \item \textbf{DRIVER\_OK:} Die Initialisierung ist abgeschlossen; das Gerät ist betriebsbereit.
\end{enumerate}

Schlägt ein obligatorischer Schritt fehl, wird der Fehlerstatus gesetzt und die
Initialisierung abgebrochen. Dadurch wird verhindert, dass ein teilweise
initialisiertes Gerät in einen undefinierten Laufzeitzustand gelangt.

\section{Virtqueue-Einrichtung und DMA-Verwaltung}
Das Kommunikationsrückgrat des Treibers ist die \textbf{Split Virtqueue}. Die
Funktion \texttt{virtqueue\_setup()} übernimmt die zentrale Speicherallokation
und Konfiguration der Queue-Strukturen.

\subsection{Speicherausrichtung und Allokation}
Da der Hypervisor (QEMU/KVM) direkt auf den RAM des Gasts zugreift, muss der
verwendete Speicher physisch zusammenhängend und passend ausgerichtet sein. Der
Treiber verwendet \texttt{os\_alloc\_dma\_pages()}, um seitenausgerichtete
Blöcke zu reservieren und sowohl virtuelle als auch physische Adressen für die
weitere Queue-Konfiguration bereitzustellen.

Innerhalb des allokierten DMA-Bereichs werden drei logisch getrennte Strukturen
angelegt:
\begin{itemize}
    \item die \textbf{Deskriptortabelle},
    \item der \textbf{Available Ring},
    \item der \textbf{Used Ring}.
\end{itemize}

Die physischen Basisadressen dieser Bereiche werden anschließend in die
Common-Configuration-Register des Geräts eingetragen. Erst danach kann die
jeweilige Virtqueue aktiviert und für den Befehlsverkehr verwendet werden.

\section{Das Control-Command-Protokoll}
Die GPU wird über Befehle gesteuert, die an die \textit{Control Queue}
(Warteschlange 0) gesendet werden. Die Funktion
\texttt{control\_queue\_send()} implementiert einen synchronen Wrapper um die
grundsätzlich asynchrone Virtqueue-Kommunikation.

Dieser Ansatz erlaubt es der höheren Initialisierungslogik, mit einem einfachen
Request-Response-Modell zu arbeiten, während die eigentliche Übertragung
weiterhin dem VirtIO-Ringmechanismus folgt.

\subsection{Vorab allozierte Befehlspuffer}
Um Fragmentierung zu vermeiden und die Befehlsübertragung reproduzierbar zu
halten, nutzt der Treiber vorab allozierte DMA-Puffer für Requests und
Responses. Diese Puffer werden während der Initialisierung der Control Queue
reserviert und anschließend für sämtliche Kommandos wiederverwendet.

Die synchrone Übertragung eines GPU-Kommandos erfolgt über genau zwei
Deskriptoren:
\begin{itemize}
    \item einen \textbf{Request-Deskriptor} für den Befehl,
    \item einen \textbf{Response-Deskriptor} für die Antwort des Geräts.
\end{itemize}

Der Request-Puffer wird vom Gerät gelesen, während der Response-Puffer vom Gerät
beschrieben wird. Beide Puffer werden über eine Deskriptorkette logisch zu
einem gemeinsamen Kommando verbunden.

\begin{center}
\begin{tikzpicture}[font=\small,
    mem/.style={rectangle, draw, fill=orange!5, text width=5.5cm, minimum height=0.8cm, text centered}]

    \node [mem] (req) {Anforderungspuffer (\texttt{virtio\_gpu\_ctrl\_hdr} + Nutzdaten)};
    \node [mem, below=0.2cm of req] (res) {Antwortpuffer (\texttt{virtio\_gpu\_ctrl\_hdr} + Fehlercode)};

    \node [left=0.5cm of req] {Gast-RAM (DMA)};
    \draw [->, thick, dashed] (req.east) -- +(1,0) node[right] {Gerät liest};
    \draw [<-, thick, dashed] (res.east) -- +(1,0) node[right] {Gerät schreibt};
\end{tikzpicture}
\captionof{figure}{Speicheraufteilung für Anforderung und Antwort der Control Queue}
\end{center}

\subsection{Signalisierung und Polling}
Nachdem Request- und Response-Puffer vorbereitet wurden, veröffentlicht der
Treiber den Kopf der Deskriptorkette im Available Ring und benachrichtigt das
Gerät über das Notify-Register. Anschließend wartet der Treiber auf den
Abschluss im Used Ring.

Der Ablauf ist dabei:
\begin{enumerate}
    \item Request in den DMA-Puffer kopieren,
    \item Response-Puffer vorbereiten beziehungsweise leeren,
    \item zwei Deskriptoren aufsetzen und verkettet publizieren,
    \item Kettenkopf in den Available Ring eintragen,
    \item Gerät per \texttt{virtqueue\_notify()} benachrichtigen,
    \item auf Completion über \texttt{virtqueue\_wait\_used()} warten,
    \item Antwort in den Aufruferpuffer übernehmen.
\end{enumerate}

Da die aktuelle Implementierung für Einrichtungs- und Kontrollkommandos ein
synchrones Verhalten benötigt, wird eine Polling-Schleife mit Watchdog-Timer
verwendet. Antwortet das Gerät nicht innerhalb des vorgegebenen Zeitfensters,
wird ein Fehler protokolliert und der betreffende Schritt als fehlgeschlagen
behandelt.

\section{Framebuffer und Ressourcenlebenszyklus}
Das eigentliche Ziel des Treibers ist die Bereitstellung eines sichtbaren
Framebuffers. Dies wird in \texttt{create\_framebuffer\_resource()} realisiert,
wo lokaler DMA-Speicher, VirtIO-GPU-Ressource und Scanout logisch miteinander
verbunden werden.

\subsection{Die 2D-Ressourcen-Pipeline}
Die Erstellung und Sichtbarmachung eines Framebuffers folgt einer klaren
Befehlskette:
\begin{enumerate}
    \item \textbf{RESOURCE\_CREATE\_2D:} Allokiert eine hostseitige 2D-Ressource für die angegebene Auflösung.
    \item \textbf{RESOURCE\_ATTACH\_BACKING:} Verbindet Gast-RAM-Seiten mit der hostseitigen Ressource.
    \item \textbf{SET\_SCANOUT:} Verbindet die Ressource mit dem virtuellen Monitor.
    \item \textbf{TRANSFER\_TO\_HOST\_2D:} Überträgt die geänderten Pixelbereiche aus dem Gast-Framebuffer an den Host.
    \item \textbf{RESOURCE\_FLUSH:} Signalisiert dem Host, die aktualisierten Bereiche sichtbar zu machen.
\end{enumerate}

Diese Trennung zwischen lokalem Schreiben in den Gast-Framebuffer und expliziter
Präsentation gegenüber dem Host ist zentral für die Effizienz des Treibers. Der
Gast kann seine Pixel zunächst vollständig lokal erzeugen, bevor nur die
tatsächlich relevanten Bereiche an das Gerät übertragen und sichtbar gemacht
werden.

\section{Dynamisches Rendering, Animationsthread und Laufzeit-Umschalter}
Nachdem die grundlegende statische Ausgabe (Schachbrettmuster) erfolgreich
validiert wurde, wurde der Treiber um eine vollständige dynamische
Rendering-Pipeline erweitert. Diese Erweiterung umfasst drei eng miteinander
verbundene Subsysteme: den Per-Frame-Renderer, den dedizierten
Animationsthread und einen Laufzeit-Moduswechsel per Tastatur.

\subsection{Per-Frame-Renderer und Dirty-Rect-Präsentation}
Der Kern der dynamischen Pipeline ist die Funktion
\texttt{virtio\_gpu\_demo\_render\_frame(uint32\_t frame\_index)} in
\texttt{virtio\_gpu\_resource.cpp}. Sie berechnet pro Frame die aktuellen
Positionen aller animierten Objekte und zeichnet die Szene in den
DMA-Framebuffer. Der \texttt{frame\_index}-Parameter steuert zeitbasierte
Animationen (Translationen, Farbverläufe) ohne externe Zeitgeber
deterministisch.

Um unnötige Übertragungen zu vermeiden, implementiert der Treiber eine
\textbf{Dirty-Rect-Präsentation}: Anstatt bei jeder Aktualisierung den
gesamten Framebuffer zu übertragen, werden nur die tatsächlich geänderten
Bereiche an den Host gemeldet.
\begin{itemize}
    \item \textbf{\texttt{present\_rect(x, y, w, h)}:} Sendet \texttt{TRANSFER\_TO\_HOST\_2D} und \texttt{RESOURCE\_FLUSH} für ein einzelnes Rechteck.
    \item \textbf{\texttt{present\_demo\_regions(...)}:} Berechnet die Vereinigungsmenge aus alter und neuer Objektposition sowie dem HUD-Bereich und ruft \texttt{present\_rect} für die betroffenen Teilbereiche auf.
\end{itemize}
Dieses Muster reduziert die pro Frame übertragene Datenmenge erheblich und
minimiert damit die Anzahl der VM-Exits, da der Hypervisor nur für die
veränderten Bildschirmregionen aktiviert wird.

\subsection{Dedizierter Animationsthread}
Da Animationen kontinuierlich im Hintergrund laufen sollen, ohne den restlichen
Kernel-Betrieb zu blockieren, wird die Render-Schleife in einem eigenen Thread
ausgeführt. Der Einstiegspunkt in \texttt{virtio\_gpu\_autostart.cpp} ist
\texttt{virtio\_gpu\_start\_demo\_animation()}, die nach dem erfolgreichen
Abschluss der Initialisierung (\texttt{virtio\_gpu\_init\_task()}) aufgerufen
wird. Die eigentliche Schleife läuft in
\texttt{virtio\_gpu\_animation\_task()}, die
\texttt{virtio\_gpu\_demo\_render\_frame()} mit einem inkrementierenden
Frame-Index aufruft.

\begin{center}
\begin{tikzpicture}[node distance=1.2cm, >=Stealth,
    block/.style={rectangle, draw, fill=blue!10, text width=5.5cm, text centered, minimum height=1cm, rounded corners}]
    \node[block, fill=green!10] (init) {\texttt{virtio\_gpu\_init\_task()}\\(Probe + Handshake + Framebuffer)};
    \node[block, below=of init] (start) {\texttt{virtio\_gpu\_start\_demo\_animation()}\\(Thread-Start)};
    \node[block, fill=orange!10, below=of start] (task) {\texttt{virtio\_gpu\_animation\_task()}\\(Render-Loop)};
    \node[block, fill=blue!5, below=of task] (render) {\texttt{virtio\_gpu\_demo\_render\_frame(frame\_idx++)}};
    \node[block, fill=green!5, below=of render] (present) {\texttt{present\_demo\_regions(...)}};
    \draw[->, thick] (init) -- (start);
    \draw[->, thick] (start) -- (task);
    \draw[->, thick] (task) -- (render);
    \draw[->, thick] (render) -- (present);
    \draw[->, thick] (present.east) -- ++(1,0) -- ++(0,2.5) -- ++(-1,0) -- (task.east);
\end{tikzpicture}
\captionof{figure}{Ablauf des dedizierten Animationsthreads}
\end{center}

\subsection{Laufzeit-Moduswechsel per Tastatur}
Eine wichtige Erweiterung ist der interaktive Wechsel zwischen dem dynamischen
Animationsmodus und dem statischen Schachbrettmodus zur Laufzeit. Die Taste
\texttt{C} (Scan Code Set 1: Make-Code \texttt{0x2e}) wird in
\texttt{Keyboard::trigger()} erkannt und löst den Aufruf von
\texttt{virtio\_gpu\_toggle\_checker\_mode()} aus. Der globale Zustand wird
durch eine atomare Variable gespeichert.

\begin{lstlisting}[language=C++, basicstyle=\small\ttfamily, frame=single, caption=Moduswechsel-Logik in \texttt{virtio\_gpu\_resource.cpp}]
std::atomic<bool> g_checker_mode{false};

void virtio_gpu_toggle_checker_mode() {
    bool old = g_checker_mode.load();
    g_checker_mode.store(!old);
    if (!old) {
        // Wechsel zu statischem Modus: Schachbrett zeichnen
        draw_checkerboard_full();
    }
}
\end{lstlisting}

Die Funktion \texttt{virtio\_gpu\_demo\_render\_frame()} wertet
\texttt{g\_checker\_mode} zu Beginn jedes Frames aus. Ist der statische Modus
aktiv, wird der dynamische Render-Pfad übersprungen und der Schachbrettinhalt
bleibt unverändert auf dem Bildschirm. Dadurch können beide Modi ohne Neustart
des Kernels direkt verglichen werden.

\subsection{Geänderte Quelldateien im Überblick}
Tabelle \ref{tab:changed_files} fasst die im Rahmen der dynamischen Erweiterung
geänderten Dateien zusammen.

\begin{table}[ht]
\centering
\caption{Im Rahmen der dynamischen Erweiterung geänderte Quelldateien}
\label{tab:changed_files}
\begin{tabular}{l p{8.5cm}}
\toprule
\textbf{Datei} & \textbf{Änderungen} \\ \midrule
\texttt{virtio\_gpu\_resource.cpp} & Dynamischer Renderer (\texttt{virtio\_gpu\_demo\_render\_frame}), Dirty-Rect-Präsentation (\texttt{present\_rect}, \texttt{present\_demo\_regions}), atomarer Modus-Toggle, \texttt{draw\_checkerboard\_full} \\
\texttt{virtio\_gpu\_autostart.cpp} & Start des Animationsthreads (\texttt{virtio\_gpu\_start\_demo\_animation}, \texttt{virtio\_gpu\_animation\_task}); der Init-Task startet die Animation nach erfolgreicher Probe \\
\texttt{virtio\_gpu.hpp} & Öffentliche API erweitert um \texttt{virtio\_gpu\_start\_demo\_animation()} und \texttt{virtio\_gpu\_toggle\_checker\_mode()} \\
\texttt{virtio\_gpu\_internal.hpp} & Interner Vertrag erweitert um \texttt{virtio\_gpu\_demo\_render\_frame(uint32\_t)} \\
\texttt{Keyboard.cpp} & Erkennung der Taste \texttt{C} (Make-Code \texttt{0x2e}) und Aufruf von \texttt{virtio\_gpu\_toggle\_checker\_mode()} \\ \bottomrule
\end{tabular}
\end{table}

\section{Implementierung der HAL-Schnittstellenschicht}
Die Hardwareabstraktionsschicht (\texttt{os\_hal\_hhu.cpp}) ist der Teil des
Treibers, der die internen Abläufe von hhuOS konkret anspricht. Sie bildet
generische Anfragen, etwa das Abbilden von MMIO-Bereichen, auf spezifische
Kernel-Dienste ab. Dieses Design folgt der Mikrokernel-Philosophie, indem es
die Hardwarezugriffslogik klar von der zentralen Protokollbehandlung trennt
\cite{Liedtke1995}.

\section{Fazit der Implementierung}
Der implementierte Treiber verbindet den hhuOS-Kernel erfolgreich mit der GPU
des Hosts. Durch die Nutzung paravirtualisierter Virtqueues und expliziter
DMA-Verwaltung erreicht der Treiber geringe Latenzen bei
Grafikaktualisierungen bei gleichzeitig geringem Speicherbedarf im
Gast-Betriebssystem. Die dynamische Erweiterung um Per-Frame-Rendering,
Dirty-Rect-Präsentation, einen dedizierten Animationsthread und einen
interaktiven Moduswechsel transformiert den Treiber von einem einfachen
Initialisierungsnachweis hin zu einem vollständigen Laufzeit-Grafikstack.

\newpage